\chapter{Discussion}
\section{Overall Shortcomings}
\subsection{Walking for Long Periods}\label{sec:WLP}
The first drawback is the weight of the footwear when fully equipped. With the main battery, energy harvesting components, and protective housing installed, the total mass reaches 700–900 grams, which can lead to leg fatigue during prolonged use. This is a significant disadvantage, especially given that E-ShoeS are intended for daily walking and running, where the substantial added weight makes extended movement more tiring.

The physical design incorporates a high-capacity battery and integrated electronic hardware, meaning the footprint of the footwear is larger than its name might suggest. While these components are essential for the intended functions, their total weight forces the user to make a trade-off. The product may offer versatile smart features, but it can become uncomfortable to wear over long periods.

\subsection{High Production Costs and Limited Accessibility}\label{sec:HPCLA}
The second drawback is cost. Materials such as carbon or titanium drive up expenses thereby limiting accessibility for the broader consumer base. This drawback reflects a tension between technical ambition and market scalability. While certain materials may be essential to realizing a product exactly as envisioned, their high cost inevitably limits the number of people who can access it. Nevertheless, this is an issue that warrants careful consideration. 

\section{Technical Limitations}
\subsection{Mechanical Fatigue of PZT Energy Harvesters}\label{sec:MFPEH}
The setup incorporates GZT-5H piezoelectric material beneath the heel and the forefoot to harvest energy from the forces generated during walking. While these locations are effective for energy collection, prolonged use can negatively impact the functionality of the materials as continuous application of force is detrimental to them. Therefore, the challenge with PZT lies not only in generating the initial electrical response but also in maintaining it continuously and without error. Achieving both simultaneously is extremely difficult. For EShoeS, this inconvenience would create a feeling of insecurity for the wearer.

\subsection{Environmental Hazards and Battery Protection}\label{sec:EHBP}
Because kinetic energy harvesting relies solely on movement, outdoor weather conditions like rain or mud can impact exterior shoe parts, requiring protective waterproofing for internal sensors and wiring.

\subsection{Risk of Battery Impact and Overheating}\label{sec:RBIO}
The energy storage system utilizes a 3.7 V, 4000 mAh Li-Po battery. Because this high-capacity battery supplies continuous power to the lighting and electronic subsystems inside a compact sole, an additional protection system is needed within the battery compartment.

According to technical specifications, severe impacts or exposure to excessive heat can lead to fire or explosion. Therefore, the placement of the battery integrated into the footwear must be carefully considered to minimize the risk of physical impact and to ensure that heat can be effectively dissipated if the battery is subjected to a significant external heat source.

\subsection{Fatigue of LED Circuitry During Sole Bending}\label{sec:FLCDSB}
The soles of the shoes are encircled by RGB LED strips. Placing the circuitry here makes it susceptible to bending during vigorous movement or sudden changes in direction, which creates a risk of the circuit breaking. While these light strips may look great when the wearer is standing still, they fail to perform effectively during movement.

This detail also has implications for the visual presentation of the EShoeS. If the circuitry cannot maintain its intended shape during vigorous movement, its integrated lighting function can no longer reliably ensure safety.

\section{Proposed Engineering Solutions}
\subsection{Composite reinforcement for the PZT}
To address the risk mentioned in Section~\ref{sec:MFPEH}, the proposed base on PZT-5H design will be reinforced with a soft composite, primarily to enable the energy harvester to withstand repetitive mechanical impact while maintaining its piezoelectric functionality.

This upgrade is intended to ensure the PZT assembly operates more smoothly over successive cycles, with the proposal specifying that it must withstand over 5 million compression cycles without compromising other functions.

\subsection{Mechanical and thermal protection for batteries}
As the limitations mentioned in Section~\ref{sec:RBIO}, the solution is utilizing a carbon shell made of Kevlar fibers combined with a thermally conductive adhesive and an Aerogel insulation layer. The casing serves to protect the battery from physical impact, while other materials perform thermal functions.

Thermal interface material and thermal insulation film serve two completely opposite functions: the former is used to facilitate heat transfer within an assembly while Aerogel insulation film acts as a thermal barrier.

Addressing this weakness significantly extends the lifespan of the footwear and enhances the user experience. However, this approach does not entirely eliminate the risk of fire or explosion.

\subsection{Flexible Printed Circuit Design for the RGB LED Strip}
To address the challenges encountered in Section~\ref{sec:FLCDSB}, the proposed solution involves using a flexible printed circuit made of rolled copper, incorporating zigzag-shaped conductive traces. This modification is better suited for repetitive movements.

Changing both the materials and the design of the conductive paths are two closely complementary tasks. Flexible circuits tend to adapt well to bends while zigzag-patterned paths are used to accommodate specific operational requirements.

The plan is to achieve over 1 million bending cycles. This is a durability target specifically for the circuit. Therefore, this change could make that goal more attainable.

\subsection{Commercial Product Segmentation}
In Section~\ref{sec:HPCLA}, cost also presents a significant challenge. The proposed solution is to divide the product line into two distinct segments “basic” and “sport.” The basic segment will utilize standard composite materials while the sport segment will incorporate high-end materials such as carbon fiber, titanium and others.

This proposal offers buyers a wider range of options, ensuring they do not have to commit a large sum of money to a product that may not suit their needs. At the same time, it allows those who are interested to purchase high-end products to experience them.

It is expected that there will be a significant price gap between the two segments, allowing users who opt for the basic version to save on costs and reduce their financial burden.
